% Generated draft: requires mathematical and source review.
\documentclass[11pt]{article}
\usepackage[margin=25mm]{geometry}
\usepackage{fontspec}
\setmainfont{Latin Modern Roman}
\usepackage{amsmath,amssymb,mathtools,unicode-math}
\setmathfont{Latin Modern Math}
\usepackage{enumitem,needspace,xurl,hyperref}
\setlength{\parindent}{0pt}
\setlength{\parskip}{5pt}
\setlength{\emergencystretch}{3em}
\newcommand{\N}{\mathbb{N}}
\newcommand{\Z}{\mathbb{Z}}
\newcommand{\Q}{\mathbb{Q}}
\newcommand{\R}{\mathbb{R}}
\newcommand{\C}{\mathbb{C}}
\newcommand{\F}{\mathbb{F}}
\newcommand{\A}{\mathbb{A}}
\newcommand{\PP}{\mathbb{P}}
\newcommand{\E}{\mathbb{E}}
\newcommand{\eps}{\varepsilon}
\newcommand{\dd}{\,\mathrm{d}}
\begin{document}
\section*{UnsolvedMath ID: 30000853}
\textbf{Source problem code:} catalogue entry
\subsection{Definitions and mathematical statement}
Let \(g\geq1\), let \(\mathcal A_g\) denote the moduli space of principally polarized abelian varieties of dimension \(g\), and let \(X_g\) be the intended toroidal compactification, called \(\mathcal A_g\) itself in the extraction. Set \(G=g(g+1)/2=\dim X_g\). Write \(\lambda=c_1(\det\mathbb E)\) for the Hodge line-bundle class, where \(\mathbb E\) is the rank-\(g\) Hodge bundle with its intended extension, and write \(D\) for the intended boundary divisor class.

The conventional interpretation of the source's Hodge/boundary intersection number is
\[
 a_N(X_g)=\deg\bigl(\lambda^{G-N}D^N\cap[X_g]\bigr),
 \qquad 0\leq N\leq G.
\]
The record does not explicitly give this formula, the initial compactification, or the boundary normalization, so they must be confirmed in the report; this interpretation is not a recovered definition from a supplied formula.

With the intended \(a_N\), the conjectured vanishing is
\[
 \forall g\geq1\ \forall N\in\{0,\ldots,G\},\qquad
 G-N\notin\left\{\frac{k(k+1)}2:0\leq k\leq g\right\}
 \ \Longrightarrow\ a_N=0.
\]
Equivalently, \(a_N\ne0\) is allowed only when \(G-N\) is the dimension of a Satake stratum \(\mathcal A_k\), since
\[
 \mathcal A_g^{\mathrm{Sat}}=\coprod_{k=0}^g\mathcal A_k,
 \qquad\dim\mathcal A_k=\frac{k(k+1)}2.
\]
The statement is a necessary condition for nonvanishing; it does not assert nonzero values at every allowed index. The case \(k=0\) includes the zero-dimensional stratum.

The second request asks whether the same vanishing pattern holds for other ``reasonable'' toroidal compactifications, with appropriately defined Hodge and boundary classes there. The class of compactifications covered by reasonable, and the reference compactification for comparison, are not specified. No equality of numerical intersection values across compactifications is requested.

The supplied assessment dated 2026-08-21 reports proofs in substantial ranges and computations, but says these do not establish the all-genera conjecture. It gives neither the exact proved ranges nor an unconditional general theorem.
\subsection{Short English statement}
The conjecture permits a Hodge/boundary top intersection to be nonzero only when the remaining Hodge exponent is a triangular number corresponding to a Satake-stratum dimension. It also asks whether this support pattern persists for other toroidal compactifications, whose class is not defined in the extraction.
\subsection{Sources}
\begin{itemize}
\item Local UnsolvedMath registry, permanent ID \(30000853\), exact source code \texttt{catalogue entry}, proposed year 2007; no individual proposer is supplied.
\item \emph{Komplexe Algebraische Geometrie (2008)}, Oberwolfach Reports, DOI 10.4171/owr/2007/47, pp. 2791--2850. \url{https://doi.org/10.4171/owr/2007/47}.
\item Supplied reference: Erdenberger--Grushevsky--Hulek, ``intersection numbers,'' \url{https://arxiv.org/abs/0707.1274}.
\item The supplied \texttt{corrected\_verified} note reports checking mathematical content against the official report. The open label and assessment checked 2026-08-21 are dated supplied triage requiring independent verification of the conjecture's status and the stated partial ranges.
\end{itemize}
\end{document}
